Write $q=n/d$ and $\sigma=s/h$, and work throughout with $2a\le s\le3$. Take $0<\gamma<w_0/8$ initially. Constants in the wall estimates can be taken independent of $\gamma$. Prefactors in the interior estimates may depend on $\gamma$. A power $\sigma^m$ below denotes a fixed power, which can be increased finitely many times. The inverse auxiliary chart
\[
Y=\frac{\sqrt2\tau+aQ}{1+W\cdot W/2},\qquad P=WY,\qquad Z=\sqrt2Y-\tau
\]
fixes the shape of the fiber in $(W,Q)$ and gives $y\asymp s$. Where the wall $r=0$ meets the support of the cutoff, $|Q|\le1/20$ forces $|W|\ge\sqrt3w_0/2$. We use a separate interior region with $|W|\le\gamma$ and $|Q|\le1/400$.

\smallskip
\noindent\textit{Spatial average and upper bounds on the current.}
For a fixed dimensional constant $D\ge1$ to be chosen, set $P_*(X)=(J_{Ds}*|g_T|^2)^{1/d}$ and $M=2n+4$. The kernel $J_{Ds}$ is strictly positive and $\|g_T\|_2=1$, so $P_*>0$. Termwise differentiation of the periodized kernel through order two gives $|\nabla_X^j J_{Ds}|\le C(Ds)^{-j}J_{Ds}$, since the rescaled Euclidean kernel has the same relative derivative bounds. Differentiating the $d$th root and using the order-one estimate for its quadratic term gives
\[
|\nabla_X^j P_*|\le C(Ds)^{-j}P_*\quad(j=1,2),\qquad
\int_{\mathbb T^n}P_*^{d+1}\,\mathrm dX\le\|g_T\|_{p_c}^{p_c}.
\]
The last inequality is Young's inequality in $L^q$ for the probability kernel, using $d+1=n$ and $2q=p_c$.

By Proposition~\ref{ms:prop-2-2}, the entries of the column whose outer product is $A^0$ are sums of scalar multiplier outputs placed in quaternionic slots, with scalar symbols a bounded smooth factor times
\[
(\kappa/y)^{d/2}e^{-\Psi},\qquad \Psi=\kappa y|W-v|^2/2.
\]
Multiply the symbols by a fixed smooth annular cutoff equal to one on the data and supported where $|v|<w_0/3$ and $\kappa>0$. In the scaled frequency variable $\eta=h\xi$, the derivatives of the symbols through order $M$ are bounded by $C(hs)^{-d/2}\sigma^m$. Indeed, $h\kappa$ and $v$ are smooth on this cutoff with uniformly bounded derivatives and $h\kappa$ is bounded away from zero. Every derivative of $\Psi$ through order $M$ costs at most $C\sigma$. Where $r=0$ and $g(Q)\ne0$, one has $|W-v|\ge\sqrt3w_0/2-w_0/3>0$ on this enlarged cap. There the derivative bound gains a factor $e^{-c_0\sigma}$, with $c_0>0$ independent of $\gamma$.

If $B$ bounds the derivatives through order $M$ of a symbol supported in the fixed scaled annulus, apply $(1-\Delta_\eta)^{n+2}$ to the symbol in the inverse Fourier integral. Its absolute Euclidean kernel is at most $CBh^{-n}(1+|X/h|^2)^{-M/2}$. Periodizing gives an absolute kernel bound by $CBJ_h$ on the torus; the same argument works in a rotated frame and its rotated lattice. By Cauchy--Schwarz against the probability kernel, the squared modulus of the output is bounded by $CB^2(J_h*|g_T|^2)$. This applies to each of the finitely many scalar slots from the current. The smoothing of $A^0$ applies one more convolution by $J_h$. On $\mathbb R^n$, either $|Y'|\ge |X|/2$ or $|X-Y'|\ge |X|/2$ in the convolution variable $Y'$, so the convolution of two copies of the positive Euclidean kernel is bounded by a constant times the kernel itself. Periodization gives $J_h*J_h\le CJ_h$. Since $Ds\ge h$, comparison term by term gives $J_h\le(Ds/h)^nJ_{Ds}$. Using $(hs)^{-d}=s^{-2d}\sigma^d$ and relabeling $m,c_0$ gives the operator norm estimates
\[
\|A_{\mathrm w}\|\le Cs^{-2d}\sigma^mP_*^d\quad(r\le0),\qquad
\|A_{\mathrm w}\|\le Cs^{-2d}\sigma^m e^{-c_0\sigma}P_*^d\quad(r=0).
\]

The time cofactor of $H(r)$ is $\det H_{\mathrm{fib}}(r)$ by Lemma~\ref{ms:lemma-3-2}; by Lemma~\ref{ms:lemma-4-1} this is comparable to $a^{-2}s^{-2d}$. A nearest lattice term in $J_{Ds}$ gives the lower bound $\inf J_{Ds}\ge c_D h^{M-n}$: the nearest representative is in a fixed bounded set, $h\le s\le3$, and the scaled kernel there is bounded below by $c_Ds^{M-n}$. Thus $P_*^d\ge c_Dh^{M-n}$. As $a\ge h$ and $50n-2\ge M-n$, this absorbs the contribution $Ch^{50n}a^{-2}s^{-2d}$ of the cofactor in the density; in particular
\[
\mu\le Cs^{-2d}\sigma^mP_*^d\quad(r\le0).
\]
We also use the weaker ceiling $\mu^{1/d}\le\lambda$ for small enough $h$ uniformly on the fibers. By Lemma~\ref{ms:wave-cap-estimates}, $\kappa\asymp h^{-1}$ and the amplitudes sum in absolute value to $O(h^{-n/2})$. Each wave weight is $O(h^{-d})$ since $y\asymp s\ge h$, and each column factor is bounded. Hence $\|(A^0)\|\lesssim h^{-n-2d}$; smoothing, cutting off and adding the preceding cofactor bound give $\mu\lesssim h^{-n-2d}+h^{50n-2-2d}$. This is at most $h^{-30nd}=\lambda^d$ for sufficiently small $h$.

\smallskip
\noindent\textit{The wall comparison and the flux.}
Put $R=H_{\mathrm{fib}}(r)$ and $B_{\mathrm{fib}}=H_{\mathrm{fib}}(B_0)$. Lemma~\ref{ms:lemma-4-1} gives $R\ge cs^{-2}I$ and $G_{\mathrm{fib}}(r)G_{\mathrm{fib}}(r)^*\le CR$. The Hessian product rule of Lemma~\ref{ms:lemma-3-1} shows that $H_{\mathrm{fib}}(rP_*)$ is the sum of $P_*R$, $rH_{\mathrm{fib}}(P_*)$, and the two mixed gradient products. Against a quaternionic vector $e$, the pairings of $G_{\mathrm{fib}}(r)$ and $G_{\mathrm{fib}}(P_*)$ with $e$ are bounded in modulus by $C(e^*Re)^{1/2}$ and $CD^{-1}P_*(e^*Re)^{1/2}$, respectively. The second derivatives of $P_*$ and $|r|\le1$ bound the cost of $rH_{\mathrm{fib}}(P_*)$ by $CD^{-2}P_*R$ in quadratic form. Hence $H_{\mathrm{fib}}(rP_*)\ge(1-C/D)P_*R$; fix $D$ large so the coefficient is at least $1/2$.

The regularizer gives $\mu>0$ by the positivity of $R$ and of $A_{\mathrm w}$. The background $B_{\mathrm{fib}}$ has rank one and positive eigenvalue comparable to $\lambda$ by Lemma~\ref{ms:lemma-4-1}. The Dirichlet theorem and comparison of Theorem~\ref{ms:thm-a-1} apply because $r$ has strictly positive fiber Hessian, $\det B_{\mathrm{fib}}=0$, and the prescribed determinant is $\lambda\mu>0$. In particular the smooth $u$ exists and $u\le0$. For $v_*=C_1\sigma^{m_1}rP_*$, compare its positive definite Hessian plus the background to $B_{\mathrm{fib}}+\tfrac12 c C_1\sigma^{m_1}P_*s^{-2}I$. By determinant monotonicity from Lemma~\ref{ms:lemma-3-2} and the background eigenvalue bound this gives
\[
\det(B_{\mathrm{fib}}+H_{\mathrm{fib}}(v_*))
\ge c\lambda(C_1\sigma^{m_1}P_*s^{-2})^d\ge\lambda\mu
\]
for $m_1,C_1$ large enough. Since the two potentials have the same zero boundary values, comparison gives $v_*\le u\le0$. Dividing by $r<0$ and extending smoothly to the boundary as in Proposition~\ref{ms:prop-4-2} gives $0\le u/r\le C_1\sigma^{m_1}P_*$.

Proposition~\ref{ms:prop-4-2} gives smoothness of $\mathcal F$ in the entire open parameter domain: the Dirichlet solution is smooth up to the moving wall and the auxiliary chart identifies the fibers with a fixed domain. The determinant in the segment integral is polynomial, so it remains smooth at $\theta=0$. Lemma~\ref{ms:lemma-4-3}, applied on each fiber with the density $\mu=A_{TT}$ and its smooth Dirichlet solution $u$, gives
\[
 0\le \frac{n-1}{n+1}\int_{\Omega_z}(-u)\mu\,\mathrm d\zeta
 \le \mathcal F(z)\le \int_{\Omega_z}(-u)\mu\,\mathrm d\zeta .
\]
Here $u\le0$ and $\mu>0$ follow from the Dirichlet construction on each fiber. The volume comparison of Lemma~\ref{ms:lemma-4-1} gives an auxiliary volume $O(a^2s^{2d})$. Using the upper bounds on $u/r$ and $\mu$ and $\int P_*^{d+1}\,\mathrm dX\le\|g_T\|_{p_c}^{p_c}$, we obtain
\[
\mathcal F(T,\tau)\le Ca^2\sigma^m\|g_T\|_{p_c}^{p_c}.
\]

For the wall flux, Proposition~\ref{ms:prop-2-2} gives $(\partial\Phi)A^0=0$. By the gradient identity in Lemma~\ref{ms:lemma-3-1}, the Hermitian current kills $G(|\Phi|^2)$, and this survives multiplication by $g(Q)$ and convolution because $\partial\Phi$ is independent of $X$. Hence only the $|W|^2/w_0^2$ part of the full $G(r)$ contributes. Lemma~\ref{ms:lemma-4-1} gives $|G(|W|^2)|\lesssim s^{-1}$, so $G(r)^*A_{\mathrm w}G(r)\le Cs^{-2}\|A_{\mathrm w}\|$. In the $(W,Q)$ chart the coarea measure $\delta(r)\,\mathrm dP\,\mathrm dZ$ has total mass at most $Ca^2s^{2d}$ by the same volume comparison and coarea on the fixed ellipsoid. Proposition~\ref{ms:prop-4-2} bounds $\Delta_z\mathcal F$ above by the wall integral of $(u/r)G(r)^*A_{\mathrm w}G(r)\delta(r)$. The wall current bound and the upper bound for $u/r$ then yield
\[
\Delta_z\mathcal F(T,\tau)\le C(a/s)^2\sigma^me^{-c_0\sigma}\|g_T\|_{p_c}^{p_c}.
\]

\smallskip
\noindent\textit{Interior recovery and lower bound.}
Take a center $(X,P,Z)$ with $|W|\le\gamma$, $|Q|\le1/400$. The radius-$h$ ball in the $3n$ real spatial slots about the center fits inside $|W|\le2\gamma$, $|Q|\le1/100$ for sufficiently large $K\gamma$ and $K$: the chart bounds $\nabla W=O(s^{-1})$, $\nabla\Phi=O(1)$ and gives changes at most $Ch/s$ and $Ch$ along such segments. These derivative bounds hold up to the edge of the enlarged patch, so a first exit would violate the chosen bounds. The cutoff is identically one on this ball.

By the time-component identity in Proposition~\ref{ms:prop-2-2}, $|\mathsf S_0g_T|^2\le A^0_{TT}$. Multiply the reciprocal of this symbol from Lemma~\ref{ms:wave-cap-estimates} by a smooth annular cutoff equal to one on the data, supported where $|v|\le2\gamma$. On the ball $|W|\le2\gamma$, so $|W-v|\le4\gamma$ on this enlarged cap, and $\Psi\le C\gamma^2\sigma$. The derivatives through order $M$ in $\eta=h\xi$ are bounded by
\[
B_{\mathrm{inv}}=C_\gamma(hs)^{d/2}\sigma^me^{C\gamma^2\sigma}.
\]
This follows by differentiation since $h\kappa$ stays bounded away from zero and all derivatives of $\Psi$ cost at most $C\sigma$. Derivatives of the small cap cutoff affect $C_\gamma$, while the constant $C$ in the exponential is independent of $\gamma$.

Apply the kernel estimate and Cauchy--Schwarz to this inverse acting on $\mathsf S_0g_T$. At the real center $X$ and any auxiliary point in the ball it gives
\[
|g_T(X)|^2\le CB_{\mathrm{inv}}^2(J_h*A^0_{TT})(X).
\]
If $|X'-X|\le h$, the two kernel values $J_h(X'-Y')$ and $J_h(X-Y')$ are uniformly comparable, by the analogous termwise statement for the Euclidean kernel. Since $\mu\ge J_h*A^0_{TT}$ where $g(Q)=1$, this proves the pointwise lower bound on the entire ball
\[
\mu\ge m_X:=c_\gamma(hs)^{-d}\sigma^{-m}e^{-C_2\gamma^2\sigma}|g_T(X)|^2.
\]

Let $\ell$ be the Euclidean distance to the center in the spatial slots. When $m_X>0$, the quadratic $v_b=\epsilon m_X^{1/d}(\ell^2-h^2)$ has zero ball-boundary value and fiber Hessian $6\epsilon m_X^{1/d}I$ by the three-slot Hessian formula. The rank-one background and $m_X^{1/d}\le \mu^{1/d}\le\lambda$ give, for a fixed small $\epsilon>0$,
\[
\det(B_{\mathrm{fib}}+H_{\mathrm{fib}}(v_b))
\le C[(\epsilon m_X^{1/d})^n+\lambda(\epsilon m_X^{1/d})^d]
\le\lambda m_X .
\]
Both Hessians are positive definite. Since $u\le0=v_b$ on the ball boundary and $\det D_u=\lambda\mu\ge\lambda m_X$, the local comparison in Theorem~\ref{ms:thm-a-1} gives $u\le v_b$ on the ball. At the center, $(-u)\mu\ge\epsilon h^2m_X^{1+1/d}$; if $m_X=0$ this follows from $u\le0$ and $\mu\ge0$.

The auxiliary volume of the centers in $(P,Z)$ is comparable to $a^2s^{2d}\gamma^{2d}$ by the Jacobian in Lemma~\ref{ms:lemma-4-1}. Since $1+1/d=q$ and $2q=p_c$, integration over the centers, followed by the lower comparison in Lemma~\ref{ms:lemma-4-3}, gives
\[
\mathcal F(T,\tau)\ge c_\gamma a^2\sigma^{-m}e^{-C_3\gamma^2\sigma}\|g_T\|_{p_c}^{p_c};
\]
here $h^2s^{2d}(hs)^{-dq}=\sigma^{n-2}$ and the constant in the exponential remains independent of $\gamma$.

Now fix $\gamma$ so that $C_3\gamma^2<c_0/2$. Choose $K$ large enough for the ball inclusions and so that $h<1/(4K)$ gives the density ceiling used above. The interior bound is positive since $\|g_T\|_2=1$. Dividing the flux estimate by it gives $C_\gamma s^{-2}\sigma^{2m}e^{-(c_0-C_3\gamma^2)\sigma}$, bounded by $Cs^{-2}e^{-c\sigma}$ after decreasing the exponential constant. At $s=4a$ and $s=2a$ the ratios $\sigma$ are respectively $4K$ and $2K$. The energy upper bound at the first radius with $g_0=f$ and the lower bound at the second radius give the stated circle estimates.
